\documentclass[11pt]{article}
\usepackage[T1]{fontenc}
\usepackage[utf8]{inputenc}
\usepackage{amsmath,amssymb,amsthm}
\usepackage[margin=1in]{geometry}
\usepackage{array}
\usepackage[colorlinks=true,linkcolor=blue,citecolor=blue]{hyperref}
\usepackage{xcolor}

\newtheorem{theorem}{Theorem}[section]
\newtheorem{lemma}[theorem]{Lemma}
\newtheorem{proposition}[theorem]{Proposition}
\newtheorem{corollary}[theorem]{Corollary}
\newtheorem{assumption}[theorem]{Assumption}
\theoremstyle{definition}
\newtheorem{definition}[theorem]{Definition}
\newtheorem{example}[theorem]{Example}
\theoremstyle{remark}
\newtheorem{remark}[theorem]{Remark}

\newcommand{\ASSUMED}{\textcolor{orange}{\textbf{[ASSUMED]}}}
\newcommand{\PROVED}{\textcolor{green!60!black}{\textbf{[PROVED]}}}
\newcommand{\IMPORTED}{\textcolor{blue}{\textbf{[IMPORTED]}}}
\newcommand{\CONJ}{\textcolor{red}{\textbf{[CONJECTURED]}}}
\newcommand{\Erec}{E_{\mathrm{rec}}}
\newcommand{\Ops}{\mathsf{Ops}^{\mathrm{phys}}_A}

\title{Typed Pipeline for Recoverability--Rate--Power Links:\\
\large A Contract Paper with a Closed Recoverability Lane and
Falsification Criteria}
\author{Lluis Eriksson}
\date{January 2026 (v1); July 2026 (v2)}

\begin{document}
\maketitle

\begin{abstract}
We present an audit-friendly \emph{logical contract} for a multi-layer
program connecting (i) static locality/Markovness, (ii) recoverability
bounds, (iii) separation-dependent dissipation rates, and (iv)
thermodynamic maintenance power. Each interface is typed with explicit
quantifiers, tagged as
\PROVED/\IMPORTED/\ASSUMED/\CONJ, and paired with falsification routes.
We do \emph{not} claim a proof of the Clay Yang--Mills mass gap; we
separate a Clay (closed-Hamiltonian) track from an operational
(open-system/maintenance) track. As a fully closed lane inside this
paper, we prove that an exponential conditional mutual information (CMI)
decay hypothesis implies exponential recoverability via an imported
Fawzi--Renner inequality, with fidelity conventions fixed explicitly.
\emph{v2 (no v1 number is changed):} the cross-reference labels are
repaired (in v1 every assumption, lemma, theorem, remark and corollary
was typeset as ``Definition~$x.y$'') and Tables~\ref{tab:audit} and
\ref{tab:fals} are re-typeset; the constant-alignment step of
Appendix~\ref{app:cfr} is executed rather than prescribed --- in the
locked conventions (squared fidelity, $\Erec = -\log F$) the
Fawzi--Renner import yields $c_{\mathrm{FR}} = 1$ exactly
(Remark~\ref{rem:cfrone}); the deferred interfaces of
Appendix~\ref{app:deferred} are now concrete series papers and are cited
as such (the Type III interface is ai.viXra:2601.0065, the
Davies/RIP-U dynamics note is ai.viXra:2601.0064, the benchmark
infrastructure lives in the series repository); the BATO-LAW upgrade
path of Appendix~\ref{app:bato} records the partial progress of
ai.viXra:2601.0031; and a verification suite instantiates everything
instantiable: the closed recoverability lane end to end on a gapped transverse-field
Ising collar ($I \approx 0.20\,e^{-1.06 w}$, $\Erec \le I$ at every
width in the generated dataset, full $1 - F \le \Erec \le
c_{\mathrm{FR}} K e^{-\alpha\epsilon}$ chain), the exact-Markov example
at machine precision, and the dephasing example with an explicit
$\kappa^{\uparrow}/\kappa^{\downarrow}$ spread of two orders of
magnitude illustrating the directionality golden rule.
\end{abstract}

\section{Positioning and related work (minimal)}

Recoverability from small conditional mutual information (CMI) was
initiated in the finite-dimensional setting by Fawzi--Renner \cite{FR},
with subsequent developments clarifying strengthened and universal
recovery maps (see e.g.\ \cite{STH} and related work). Exponential CMI
decay in Gibbs or quasi-local settings has been studied in several
regimes in the many-body/QIT literature; in this paper we keep that
input as an explicit falsifiable hypothesis (A-CMI) rather than
embedding it as a black-box ``Gibbs theorem'' claim. The present
contribution is the \emph{typed contract} structure and the explicit
separation between proved implications and bridge assumptions
(especially toward operational power costs).

\section{Executive summary (audit view)}

\textbf{Scope and non-claims}
\begin{itemize}
\item \textbf{Not a Clay proof:} no claim is made about existence of
continuum Yang--Mills theory or its Hamiltonian spectral gap.
\item \textbf{This is a contract paper:} the value is explicit typing,
status tagging, and falsification routes.
\item \textbf{Closed lane included:} we include one fully
self-contained implication: CMI decay (assumption) $+$ FR-type
recoverability (imported) $\Rightarrow$ recoverability decay (proved).
\end{itemize}

\textbf{Status table (what is actually proved inside this PDF)}

\begin{table}[h!]
\centering
\begin{tabular}{p{3.4cm} p{2.4cm} p{3.6cm} p{4.4cm}}
\hline
\textbf{Item} & \textbf{Status} & \textbf{Depends on} &
\textbf{Falsification route} \\
\hline
L1: $1 - F \le -\log F$ & \PROVED & elementary calculus & direct check
\\
T2: FR-type inequality (constant-fixed) & \IMPORTED & finite-dimensional
QIT & disproof or convention mismatch \\
A-CMI: CMI decay hypothesis & \ASSUMED & chosen model/geometry & exhibit
family violating stated decay \\
T3: CMI decay $\Rightarrow$ recoverability decay & \PROVED & A-CMI $+$
T2 $+$ L1 & refute A-CMI or mismatch hypotheses \\
BATO-LAW: maintenance inequality & \ASSUMED & operational model
definition & violate hypothesis or give counter-protocol \\
\hline
\end{tabular}
\caption{Audit status for this paper. Anything tagged \ASSUMED\ is a
typed hypothesis (not a hidden dependency). (Re-typeset in v2; v1's
table had overlapping column text.)}
\label{tab:audit}
\end{table}

\section{Typing, notation, and objects}

\begin{remark}[Parameter convention]
\label{rem:params}
Throughout: $\epsilon$ is a \textbf{geometric} separation/collar
parameter, while $\delta$ denotes \textbf{tolerances} (regularization
thresholds, error budgets, floors). They are never interchanged.
\end{remark}

\begin{definition}[D1: $\Delta$-track coherence]
\label{def:cdelta}
Fix a dephasing (pinching) CPTP map $\Delta$. Define $C_\Delta(\rho) :=
S(\rho \| \Delta[\rho])$, where $S(\cdot\|\cdot)$ is quantum relative
entropy (with the usual support convention).
\end{definition}

\begin{definition}[D2: instantaneous coherence loss under fixed
dynamics]
\label{def:loss}
Let $T_t = e^{t\mathcal{L}}$ be a fixed CPTP Markov semigroup on states.
Define
\[
\dot{C}_{\Delta,\mathrm{loss}}(\rho) := -\left.\frac{d}{dt}\right|_{t=0}
C_\Delta(T_t(\rho)),
\]
whenever the derivative exists.
\end{definition}

\begin{remark}[A minimal sufficient condition for existence]
\label{rem:exist}
A sufficient (not necessary) condition for
$\dot{C}_{\Delta,\mathrm{loss}}(\rho)$ to exist is differentiability at
$t = 0$ of $t \mapsto C_\Delta(T_t(\rho))$. In finite dimension this
holds under mild spectral regularity (e.g.\ full-rank conditions), but
this paper does not fix a universal regularity class; we treat existence
as part of the typing of any application.
\end{remark}

\begin{definition}[D3: physically implementable local operations $\Ops$]
\label{def:ops}
Fix a reference state $\rho_0$, a tripartition $A$--$B$--$C$, and a
separation notion $\mathrm{sep}(A, C)$. Let $\Ops$ denote a declared
class of \emph{physically implementable} operations on $A$ (CPTP maps
with local ancillas allowed and discarded), realizable within stated
resource tolerances $\delta$ (time, energy, ancilla size, etc.).
\end{definition}

\begin{definition}[D4: the family $\mathcal{F}_\epsilon$ (minimal; typed
by $\Ops$)]
\label{def:family}
Define
$\mathcal{F}_\epsilon := \bigl\{\rho = \Lambda_A(\rho_0) \,:\, \Lambda_A
\in \Ops,\ \mathrm{sep}(A, C) \ge \epsilon,\ C_\Delta(\rho) > 0\bigr\}$.
\end{definition}

\begin{assumption}[REG: regularity restrictions (invoked only when
needed)]
\label{ass:reg}
When needed to exclude trivial degeneracies (e.g.\ near-fixed points or
nearly singular diagonals) we impose explicit regularity restrictions on
$\mathcal{F}_\epsilon$, such as $\mathrm{dist}(\rho, \mathrm{Fix}(T))
\ge \delta_{\mathrm{fix}}$ and/or $\Delta[\rho] \succeq
p_{\min}\mathbf{1}$. These restrictions are not part of the minimal
definition of $\mathcal{F}_\epsilon$; they are invoked only where
required and always stated at the point of use.
\end{assumption}

\begin{definition}[D5: envelopes (quantifiers explicit)]
\label{def:envelopes}
Define the domain of the ratio at scale $\epsilon$ as
$\mathcal{D}_\epsilon := \bigl\{\rho \in \mathcal{F}_\epsilon :
C_\Delta(\rho) > 0,\ \dot{C}_{\Delta,\mathrm{loss}}(\rho)\
\text{exists}\bigr\}$. Then define
\[
\kappa^{\uparrow}(\epsilon) := \sup\left\{
\frac{\dot{C}_{\Delta,\mathrm{loss}}(\rho)}{C_\Delta(\rho)} : \rho \in
\mathcal{D}_\epsilon \right\}, \qquad
\kappa^{\downarrow}(\epsilon) := \inf\left\{
\frac{\dot{C}_{\Delta,\mathrm{loss}}(\rho)}{C_\Delta(\rho)} : \rho \in
\mathcal{D}_\epsilon \right\}.
\]
\end{definition}

\begin{remark}[Empty-domain convention]
\label{rem:empty}
In any concrete application we assume $\mathcal{D}_\epsilon \ne
\emptyset$ in the regime of interest. If $\mathcal{D}_\epsilon =
\emptyset$, we adopt the convention $\kappa^{\uparrow}(\epsilon) := 0$,
$\kappa^{\downarrow}(\epsilon) := +\infty$, so that envelope statements
become vacuous rather than ambiguous. (Any alternative convention should
be stated explicitly if used.)
\end{remark}

\begin{remark}[Directionality (golden rule)]
\label{rem:golden}
A \textbf{minimum-power} statement requires a \textbf{lower envelope}
bound (involving $\kappa^{\downarrow}$), \emph{not} an upper envelope
bound (involving $\kappa^{\uparrow}$).
\end{remark}

\section{Closed recoverability lane: CMI decay $\Rightarrow$ exponential
recovery}

\begin{definition}[Fidelity convention]
\label{def:fidelity}
For density matrices $\rho, \sigma$ on the same finite-dimensional
Hilbert space, we use the \emph{squared} Uhlmann fidelity $F(\rho,
\sigma) := \|\sqrt{\rho}\sqrt{\sigma}\|_1^2 \in [0, 1]$.
\end{definition}

\begin{definition}[Recovery error]
\label{def:erec}
For a tripartite finite-dimensional state $\rho_{ABC}$ and a recovery
map $\mathcal{R}_{B\to BC}$, define
$\Erec(\rho_{ABC}; \mathcal{R}) := -\log F(\rho_{ABC}, (\mathrm{id}_A
\otimes \mathcal{R}_{B\to BC})(\rho_{AB}))$.
\end{definition}

\begin{lemma}[L1: elementary inequality]
\label{lem:l1}
\PROVED. For $F \in (0, 1]$ one has $1 - F \le -\log F$.
\end{lemma}

\begin{proof}
Define $g(F) := -\log F - (1 - F)$ on $(0, 1]$. Then $g'(F) =
-\frac{1}{F} + 1 \le 0$ for $F \in (0, 1]$ and $g(1) = 0$, so $g(F) \ge
0$ for all $F \in (0, 1]$.
\end{proof}

\begin{theorem}[T2: FR-type recoverability inequality (constant
explicit)]
\label{thm:t2}
\IMPORTED\ (finite-dimensional recoverability literature; see e.g.\
\cite{FR, STH}). There exists a universal constant $c_{\mathrm{FR}} > 0$
(depending only on the choice of fidelity/error convention) such that
for any finite-dimensional $\rho_{ABC}$, there exists a CPTP recovery
map $\mathcal{R}_{B\to BC}$ with
\[
\Erec(\rho_{ABC}; \mathcal{R}) \le c_{\mathrm{FR}}\, I(A : C | B)_\rho.
\]
\end{theorem}

\begin{remark}[Convention lock-in]
\label{rem:lockin}
This paper fixes: (i) squared fidelity (Definition~\ref{def:fidelity}),
and (ii) error functional $\Erec = -\log F$
(Definition~\ref{def:erec}). The value of $c_{\mathrm{FR}}$ should be
aligned to the imported theorem under these conventions
(Appendix~\ref{app:cfr} gives the alignment step; \emph{executed in v2},
Remark~\ref{rem:cfrone}). Using $c_{\mathrm{FR}}$ avoids silent factor
drift across conventions.
\end{remark}

\begin{assumption}[A-CMI: exponential CMI decay in a shielded geometry]
\label{ass:acmi}
\ASSUMED\ (model/geometry hypothesis). Fix a family of tripartitions
$A$--$B(\epsilon)$--$C$ (``collar'' of width/separation $\epsilon$) and
a family of finite-dimensional states $\rho_{ABC}(\epsilon)$. Assume
there exist constants $K > 0$ and $\alpha > 0$ (independent of
$\epsilon$ in the regime of interest) such that
$I(A : C | B)_{\rho(\epsilon)} \le K e^{-\alpha\epsilon}$.
\end{assumption}

\begin{remark}[Where A-CMI is expected/known]
\label{rem:acmiwhere}
A-CMI is intended to capture ``approximate quantum Markov'' behavior in
geometries where $B(\epsilon)$ shields $A$ from $C$. In concrete
applications one either: (i) imports a theorem establishing such decay
in a stated regime (e.g.\ a temperature/interaction range regime), or
(ii) tests it numerically/empirically in a benchmark family. This paper
keeps the hypothesis explicit so that failures are interpretable as
falsification rather than hidden assumption.
\end{remark}

\begin{theorem}[T3: exponential recoverability under exponential CMI
decay]
\label{thm:t3}
\PROVED\ (given Theorem~\ref{thm:t2} and Assumption~\ref{ass:acmi}).
Under Assumption~\ref{ass:acmi}, for each $\epsilon$ there exists a CPTP
recovery map $\mathcal{R}_{B\to BC}$ such that
\[
\Erec(\rho(\epsilon); \mathcal{R}) \le c_{\mathrm{FR}}\, K
e^{-\alpha\epsilon}.
\]
Moreover, $1 - F(\rho_{ABC}(\epsilon), (\mathrm{id}_A \otimes
\mathcal{R}_{B\to BC})(\rho_{AB}(\epsilon))) \le c_{\mathrm{FR}} K
e^{-\alpha\epsilon}$.
\end{theorem}

\begin{proof}
Combine Theorem~\ref{thm:t2} with Assumption~\ref{ass:acmi}. Then apply
Lemma~\ref{lem:l1}.
\end{proof}

\begin{remark}[Why T3 is still worth stating]
\label{rem:whyt3}
T3 is logically simple, but it is \emph{audit-critical}: it fixes
conventions, isolates the single falsifiable hypothesis A-CMI, and
cleanly separates the imported recoverability mechanism from downstream
``rate--power'' bridges.
\end{remark}

\begin{remark}[v2: the alignment step executed --- $c_{\mathrm{FR}} =
1$]
\label{rem:cfrone}
\PROVED\ (v2). In the locked conventions of
Definitions~\ref{def:fidelity} and \ref{def:erec}, the Fawzi--Renner
theorem \cite{FR} states $I(A : C | B) \ge -\log F$ with $F$ the squared
fidelity of the recovered state for a rotated Petz map, i.e.\ FR imports
as Theorem~\ref{thm:t2} with $c_{\mathrm{FR}} = 1$ exactly. With the
\emph{unsquared} (root) fidelity the same theorem reads $I \ge -2\log
f$, so silently mixing conventions produces a factor-$2$ drift --- the
recurring factor corrected across this series (2512.0101, 2601.0020,
2601.0034, 2601.0035, all v2). The verification suite checks the
identity $-\log F = -2\log\sqrt{F}$ and instantiates $c_{\mathrm{FR}} =
1$ on a benchmark family (gapped transverse-field Ising collar):
$\Erec \le I(A : C | B)$ at every collar width in the generated dataset,
with $I \approx 0.20\, e^{-1.06 w}$, closing the T3 chain end to end
including the $1 - F$ form via Lemma~\ref{lem:l1}. This settles the
Type I anchor; the same alignment in the split-regularized Type III
setting is Lemma 10.2 of the companion interface note
ai.viXra:2601.0065 (v2).
\end{remark}

\section{Operational power interface (typed, not proved here)}

\begin{definition}[D6: maintenance strategies and average power]
\label{def:strat}
Fix an operational control model (battery, allowed controls, coupling to
noise), and an admissible strategy class $\mathsf{Strat}$. Fix an error
budget $\delta$ and a noise semigroup $T_t$. For a target state $\rho$,
let $P(\rho; \mathsf{S}) \in [0, \infty]$ denote the average battery
power required by a strategy $\mathsf{S} \in \mathsf{Strat}$ to maintain
$\rho$ under $T_t$ within tolerance $\delta$ (as defined by that model).
\end{definition}

\begin{definition}[D7: extra maintenance power]
\label{def:pextra}
Define the extra power as
\[
P_{\mathrm{extra}}(\rho) := \max\Bigl\{0,\ \inf_{\mathsf{S} \in
\mathsf{Strat}} P(\rho; \mathsf{S}) - \inf_{\mathsf{S} \in
\mathsf{Strat}} P(\Delta[\rho]; \mathsf{S})\Bigr\}.
\]
\end{definition}

\begin{assumption}[BATO-LAW: maintenance inequality (operational
postulate)]
\label{ass:bato}
\ASSUMED\ (typed law to be proved or imported in a companion note).
Under the hypotheses of the chosen control model, for all target states
$\rho$ one has $P_{\mathrm{extra}}(\rho) \ge k_B T\,
\dot{C}_{\Delta,\mathrm{loss}}(\rho)$, whenever
$\dot{C}_{\Delta,\mathrm{loss}}(\rho)$ exists.
\end{assumption}

\begin{corollary}[C1: envelope consequence (directionality explicit)]
\label{cor:c1}
Assume Assumption~\ref{ass:bato}. Assume Assumption~\ref{ass:reg}
whenever needed to exclude degeneracies. Then for $\rho \in
\mathcal{F}_\epsilon$,
\[
P_{\mathrm{extra}}(\rho) \ge k_B T\, \kappa^{\downarrow}(\epsilon)\,
C_\Delta(\rho).
\]
\end{corollary}

\begin{proof}
By Assumption~\ref{ass:bato}, $P_{\mathrm{extra}}(\rho) \ge k_B T
\dot{C}_{\Delta,\mathrm{loss}}(\rho)$. By
Definition~\ref{def:envelopes}, $\dot{C}_{\Delta,\mathrm{loss}}(\rho)
\ge \kappa^{\downarrow}(\epsilon) C_\Delta(\rho)$ for $\rho \in
\mathcal{F}_\epsilon$.
\end{proof}

\section{Two small examples (toy but honest)}

\begin{example}[Exact Markov case ($I(A:C|B) = 0$ implies perfect
recovery)]
\label{ex:markov}
If a state satisfies $I(A : C | B)_\rho = 0$ (i.e.\ it is a quantum
Markov chain in the exact sense), then recoverability can be perfect:
there exists a CPTP recovery map $\mathcal{R}$ such that
$F(\rho_{ABC}, (\mathrm{id}_A \otimes \mathcal{R})(\rho_{AB})) = 1$,
hence $\Erec = 0$. This anchors the interpretation of T2/T3: CMI is the
static ``Markovness defect'' that controls recoverability; the
implication $I(A : C | B)_\rho = 0 \Rightarrow$ exact recovery is
standard (see e.g.\ \cite{HJPW}).
\end{example}

\begin{example}[A semigroup where $\dot{C}_{\Delta,\mathrm{loss}}$ is
directly computable (envelope intuition)]
\label{ex:dephasing}
Let $\Delta$ be dephasing in a fixed basis and let $T_t$ be pure
dephasing in the same basis (so $T_t$ commutes with $\Delta$ and damps
off-diagonal terms exponentially). In such cases, $C_\Delta(T_t(\rho))$
is typically monotone decreasing in $t$ on broad families, and
$\dot{C}_{\Delta,\mathrm{loss}}(\rho)$ exists and is computable from the
initial damping rate. This illustrates what $\kappa^{\uparrow}(\epsilon)$
and $\kappa^{\downarrow}(\epsilon)$ measure: worst-case and best-case
\emph{initial relative} coherence-loss rates over a typed family
$\mathcal{F}_\epsilon$. (We do not claim a universal closed form for
$\dot{C}_{\Delta,\mathrm{loss}}$ in this paper; the point is operational
meaning.)
\end{example}

\begin{remark}[v2: both examples instantiated numerically]
\label{rem:examples}
The verification suite instantiates Example~\ref{ex:markov} with a
random state of Hayden--Jozsa--Petz--Winter structure $\rho = \rho_{A
B_L} \otimes \rho_{B_R C}$: $I(A : C | B) = 4\times10^{-16}$ and
Petz-type reconstruction with $1 - F = 1.3\times10^{-15}$. It
instantiates Example~\ref{ex:dephasing} with a $d = 3$ pure-dephasing
semigroup with strongly unequal pair rates: $C_\Delta(T_t\rho)$ is
monotone on all sampled trajectories, $\dot{C}_{\Delta,\mathrm{loss}} >
0$ on the family, and the empirical envelopes spread by
$\kappa^{\uparrow}/\kappa^{\downarrow} \approx 100$ --- a concrete
reminder of the golden rule (Remark~\ref{rem:golden}): an upper envelope
never certifies a floor. The same directionality is exhibited
dynamically in the Davies setting by the companion ai.viXra:2601.0064
(v2).
\end{remark}

\section{Falsification matrix (minimal)}

\begin{table}[h!]
\centering
\begin{tabular}{p{3.2cm} p{2.4cm} p{4.0cm} p{4.2cm}}
\hline
\textbf{Item} & \textbf{Status} & \textbf{Verification proxy} &
\textbf{Concrete falsifier} \\
\hline
A-CMI (CMI decay) & \ASSUMED & compute/estimate $I(A:C|B)$ vs $\epsilon$
& show non-decay or slower decay than claimed \\
T2 (FR-type) & \IMPORTED & align $c_{\mathrm{FR}}$ to cited convention
(done in v2: $c_{\mathrm{FR}} = 1$) & detect convention mismatch / wrong
constant \\
BATO-LAW & \ASSUMED & explicit model calculation & explicit protocol
violating inequality \\
Lower-envelope need & \PROVED & logic in Remark~\ref{rem:golden} and
Corollary~\ref{cor:c1} & N/A (directionality statement) \\
\hline
\end{tabular}
\caption{Minimal falsification matrix for this contract paper
(re-typeset in v2).}
\label{tab:fals}
\end{table}

\section{Clay vs operational track (explicit separation)}

\textbf{Clay track (closed Hamiltonian).} The Clay mass gap concerns a
\emph{closed} Yang--Mills QFT in the continuum and its spectral gap in
an axiomatic setting.

\textbf{Operational track (this paper).} This paper treats an
open-system/maintenance track: typed inequalities relating coherence
loss and power. Any bridge from operational rate floors to closed
Hamiltonian gaps would require additional hypotheses not provided here.

\section*{Note added in v2}
No v1 definition, assumption, theorem statement or numbering is changed;
new v2 material is appended at the end of the affected sections. v2:
(1) repairs the cross-reference labels --- in v1 every reference was
typeset as ``Definition~$x.y$'' regardless of environment, so e.g.\
``given Definitions 4.4 and 4.6'' in v1 denotes Theorem~\ref{thm:t2} and
Assumption~\ref{ass:acmi}, ``apply Definition 4.3'' denotes
Lemma~\ref{lem:l1}, and ``Assume Definition 5.3'' denotes
Assumption~\ref{ass:bato}; (2) re-typesets Tables~\ref{tab:audit} and
\ref{tab:fals} (overlapping column text in v1); (3) executes the
$c_{\mathrm{FR}}$ alignment of Appendix~\ref{app:cfr}: in the locked
conventions the Fawzi--Renner import gives $c_{\mathrm{FR}} = 1$ exactly
(Remark~\ref{rem:cfrone}); (4) updates Appendix~\ref{app:deferred}: the
three deferred interfaces announced in v1 now exist as series papers and
are cited concretely; (5) records in Appendix~\ref{app:bato} the partial
progress toward BATO-LAW in ai.viXra:2601.0031 (v2); (6) adds series
positioning (v1 cited no companion papers) and a verification suite
(\texttt{verification/2601-0066/} in the series repository)
instantiating the closed recoverability lane end to end and both examples
(Remarks~\ref{rem:cfrone} and \ref{rem:examples}).

\appendix

\section{Imported theorem conventions (how to set $c_{\mathrm{FR}}$)}
\label{app:cfr}

To make Theorem~\ref{thm:t2} fully audit-aligned, pick a single source
theorem statement and restate it verbatim under the conventions: squared
fidelity $F$ (Definition~\ref{def:fidelity}), error functional $\Erec =
-\log F$ (Definition~\ref{def:erec}). Then set $c_{\mathrm{FR}}$ to the
numerical constant appearing in that statement after translating
conventions. Using $c_{\mathrm{FR}}$ prevents silent factor drift across
papers. \emph{v2: this step is executed in Remark~\ref{rem:cfrone} with
source \cite{FR}: $c_{\mathrm{FR}} = 1$ in the locked conventions
(equivalently, $2$ if the root-fidelity convention were used).}

\section{Deferred interfaces (updated in v2: now concrete)}
\label{app:deferred}

To keep this paper self-contained and citable, we do not develop here:
\begin{itemize}
\item Type III/AQFT split-regularized CMI definitions and conditional
expectations: \emph{now} the interface note ai.viXra:2601.0065 (v2),
with the CE-pullback direction discipline of ai.viXra:2601.0046 (v2);
\item Davies/Dirichlet-form identities and RIP-U envelope bounds:
\emph{now} the dynamics note ai.viXra:2601.0064 (v2), with the
Bohr-channel decomposition of ai.viXra:2601.0023 (v2);
\item non-abelian benchmark implementation contracts and reproducibility
harness: \emph{now} the benchmark notes ai.viXra:2601.0043, 2601.0044,
2601.0047, 2601.0050, 2601.0051 (all v2) with suites in the series
repository.
\end{itemize}
Nothing in Sections 4--6 depends on these components; the citations
record where each deferred interface is developed.

\section{Upgrade path for BATO-LAW (audit guidance)}
\label{app:bato}

In this paper, BATO-LAW is treated as \ASSUMED, since its validity is
operational-model dependent. To upgrade Assumption~\ref{ass:bato} to
\IMPORTED\ or \PROVED\ within a specified control model, it suffices to
provide:
\begin{itemize}
\item an axiomatized specification of the control model (battery,
admissible controls, error metric, and the maintenance criterion within
tolerance $\delta$);
\item either a complete proof within that model or an imported theorem
with hypotheses restated verbatim under the present conventions;
\item an explicit identification (or inequality) relating the model's
dissipation functional to $\dot{C}_{\Delta,\mathrm{loss}}(\rho)$ as
defined in Definition~\ref{def:loss}.
\end{itemize}
\emph{v2 note:} partial progress exists in the series:
ai.viXra:2601.0031 (v2) proves a work-cost inequality in a GNS/KMS
battery model (with the v1 sign error corrected there to a battery
free-energy \emph{decrease}), which supplies the third bullet's
dissipation-functional identification in that model class; a full
BATO-LAW derivation remains open.

\begin{thebibliography}{12}
\bibitem{FR} O. Fawzi and R. Renner, \emph{Quantum conditional mutual
information and approximate Markov chains}, Commun. Math. Phys.
\textbf{340}, 575--611 (2015).
\bibitem{STH} D. Sutter, M. Tomamichel, and P. Harremo\"es,
\emph{Strengthened monotonicity of relative entropy via pinched Petz
recovery map}, IEEE Trans. Inf. Theory \textbf{62}, 2907--2913 (2016).
\bibitem{HJPW} P. Hayden, R. Jozsa, D. Petz, and A. Winter,
\emph{Structure of states which satisfy strong subadditivity of quantum
entropy with equality}, Commun. Math. Phys. \textbf{246}, 359--374
(2004).
\bibitem{S0065} L. Eriksson, ai.viXra:2601.0065 (2026; v2 2026).
\bibitem{S0064} L. Eriksson, ai.viXra:2601.0064 (2026; v2 2026).
\bibitem{S0046} L. Eriksson, ai.viXra:2601.0046 (2026; v2 2026).
\bibitem{S0023} L. Eriksson, ai.viXra:2601.0023 (2026; v2 2026).
\bibitem{S0031} L. Eriksson, ai.viXra:2601.0031 (2026; v2 2026).
\bibitem{S0035} L. Eriksson, ai.viXra:2601.0035 (2026; v2 2026).
\bibitem{S0043} L. Eriksson, ai.viXra:2601.0043 (2026; v2 2026).
\bibitem{S0044} L. Eriksson, ai.viXra:2601.0044 (2026; v2 2026).
\end{thebibliography}

\end{document}
